\section{The Base Specification}
\label{app:base}

This appendix collects the exact specification of the base, the parts the body states once and uses throughout. The committed collision, the nine-state pair table, is given in the body where the knit is introduced, so here we record its inverse, the variant collisions that were considered and set aside, the streaming map in full, and the layout of the twenty-four directions.

Because every collision is a permutation of nine states, it has an exact inverse, and running the inverse undoes a beat. The inert and hop states are their own inverses, and only the create-flip-annihilate cycle reverses, running $(\,0,0) \to (-1,+1) \to (+1,-1) \to (0,0)$ backward. The whole beat is inverted by un-streaming, sending each tone back the way it came, and then applying the inverse collision, which is how the closed bulk runs time backward with no arrow.

The committed pair table was fixed by comparison against alternatives, each a reversible, charge-conserving collision built to test one question, and recording them keeps the one genuine open choice precise.

\begin{table}[ht]
\centering
\small
\begin{tabular}{@{}lp{0.40\textwidth}p{0.30\textwidth}@{}}
\toprule
collision & what it does & why not committed \\
\midrule
pass-through & no interaction, every tone streams forever & the trivial control \\
the pair table & the create move, plus a hop, plus inert states, pins a charge in place & committed, but it seals the radiation channel \\
momentum-conserving rotate & moves a head-on pair between lines, leaves a lone charge ballistic & radiates but has no create move, no arrow \\
hop-off pair table & keeps the create cycle, lets lone charges stream & the create destabilizes the vacuum globally \\
leaky confine & confines a charged line, streams a neutral one & no create move, the vacuum stays empty \\
sticky reflect & reflects a crowded dock, leaves a lone tone & elastic, it does not capture \\
\bottomrule
\end{tabular}
\caption{The committed collision against the alternatives. The one open choice is which collision a self uses, the pinning pair table with the arrow against the radiating momentum-conserving rule without it.}
\label{tab:variants}
\end{table}

Streaming is a pure relabeling. After the collision, the tone now occupying a site moves to the neighbouring dock that the site points at, with its value unchanged, so the value in a given site of a given dock after streaming is the value that sat, before streaming, in that same direction at the neighbour on the opposite side. The map is a bijection on the whole mesh and is exactly invertible, and a free tone therefore travels one dock per beat in a straight line, its momentum being simply which site it sits in.

The twenty-four directions are the $D_4$ roots, the vectors with two entries $\pm 1$ and two entries zero. They fall into six axis-pairs, each pair carrying the four sign choices, $6 \times 4 = 24$, and every direction has a unique opposite, its negation, which is what defines the twelve lines the collision runs on. Composing two directions is quaternion multiplication of the corresponding unit Hurwitz quaternions, which stays inside the twenty-four, so the directions are a closed finite group, the binary tetrahedral group, with symmetry group $F_4$ of order $1152$. A dock's entire state is the ordered list of twenty-four ternary values on these directions, and the state of the universe is the tones of all docks at one beat, an array of trits and nothing more.
